\usepackage[scr=rsfs,cal=euler]{mathalfa}
\multlinegap0pt

\hyphenation{assu-me assump-tion assump-tions Axiom being coho-mol-o-gy deduce defined denote dynam-i-cal extend Extend-ing hyper-bol-ic intrin-sic order sequen-ce}

\newcounter{toto}
\renewcommand{\thetoto}{\alph{toto}}

\newenvironment{enumeratei}
{\bgroup\def\theenumi{\roman{enumi}}\def\theenumii{\arabic{enumii}}\begin{enumerate}}
{\end{enumerate}\egroup}

\newcommand\mto{\mathchoice{\longmapsto}{\mapsto}{\mapsto}{\mapsto}}

\let\oldbigsqcup\bigsqcup
\renewcommand{\bigsqcup}{\mathop{\textstyle\oldbigsqcup}\displaylimits}

\newcommand{\Changel}{\mathcode`l="7160}
\newcommand{\Changelback}{\mathcode`l="716C}
\newcommand{\NoChangel}[1]{%
\expandafter\let\csname old\string#1\endcsname=#1
\let#1=\relax
\newcommand{#1}{\mathop{\mathcode`l="716C\csname old\string#1\endcsname\mathcode`l="7160}\displaylimits}%
}

\NoChangel{\log}
\NoChangel{\ln}
\NoChangel{\lim}
\NoChangel{\limsup}
\NoChangel{\liminf}
\NoChangel{\varlimsup}
\NoChangel{\varliminf}
\NoChangel{\varinjlim}
\NoChangel{\varprojlim}

\newcommand{\lcr}{[\![}
\newcommand{\rcr}{]\!]}

\DeclareMathOperator{\Reel}{Re}
\newcommand{\loc}{\mathrm{loc}}
\newcommand{\ro}{\mathrm{o}}
\newcommand{\rs}{\mathrm{s}}
\newcommand{\rso}{\mathrm{so}}
\newcommand{\ru}{\mathrm{u}}
\newcommand{\ruo}{\mathrm{uo}}
\newcommand{\KS}{\mathrm{KS}}
\newcommand{\rSS}{\mathrm{SS}}
\newcommand{\ST}{\mathrm{ST}}
\newcommand{\sym}{\mathrm{sym}}

\newcommand{\sep}{\::\:}
\let\dpl\displaystyle
\newcommand{\Psfrac}[2]{(\sfrac{#1}{#2})}

\graphicspath{{meddane_figures}}

\usepackage{tikz}
\usetikzlibrary{arrows}
\usetikzlibrary[patterns]
\usetikzlibrary{matrix}
\usetikzlibrary{cd}
\tikzcdset{arrow style=Latin Modern}
\newcommand{\dradpl}{{\hspace*{-2.5mm}\begin{tikzcd}[ampersand replacement=\&,column sep = scriptsize]\mbox{}\ar[r,dashed]\&\mbox{}\end{tikzcd}\hspace*{-2.5mm}}}
\newcommand{\dratst}{{\hspace*{-2.5mm}\begin{tikzcd}[ampersand replacement=\&,column sep = small]\mbox{}\ar[r,dashed]\&\mbox{}\end{tikzcd}\hspace*{-2mm}}}
\let\drascst\dratst
\newcommand{\dra}{\mathchoice{\dradpl}{\dratst}{\drascst}{\drascst}}

%\usepackage{enumitem}
%\usepackage[hypcap=false]{caption}
%\usepackage[all]{xy}
%\usepackage{url}
%\usepackage{xspace}
%\usepackage{subcaption}
%\captionsetup{labelfont=rm,textfont=normalfont}
%\DeclareCaptionSubType*{figure}
%\renewcommand\thesubfigure{\thefigure(\alph{subfigure})}
%\usepackage{textcomp}
\usepackage{appendix}
%\usepackage[dvips]{graphicx}
%\usepackage{pgf}
\usepackage{pdfpages}
%\usepackage[left = 3cm, right = 3cm]{geometry}
%\usepackage{hyperref}
%\hypersetup{
%colorlinks=true, %colorise les liens
%breaklinks=true, %permet le retour à la ligne dans les liens trop longs
%urlcolor= blue, %couleur des hyperliens
%linkcolor= blue, %couleur des liens internes
%citecolor=blue, %couleur des références
%}
\usepackage{faktor}

%\overfullrule=1mm

%% Réglages personnalisés
%\textheight = 21cm
%\textwidth = 13cm
%
%% Centrer le texte horizontalement
%\oddsidemargin = 41.5pt % LaTeX accepte aussi en points
%\evensidemargin = 41.5pt % utile si document recto-verso
%
%\raggedbottom


\theoremstyle{definition}
\newtheorem{definition}{Definition}[section]
\newtheorem{definitions}[definition]{Definitions}
\newtheorem{defprop}[definition]{Definition/Proposition}
\newtheorem{example}[definition]{Example}
\newtheorem{notation}[definition]{Notation}
\newtheorem{remark}[definition]{Remark}

\theoremstyle{plain}
\newtheorem{thm}{Theorem}
\newtheorem{prop}[definition]{Proposition}
\newtheorem{lemma}[definition]{Lemma}
\newtheorem{sublemma}[definition]{Sub-Lemma}
\newtheorem{corollary}[definition]{Corollary}

\newcommand{\Acal}{\mathcal{A}}
\newcommand{\A}{\emph{A}}
\newcommand{\Ccal}{\mathcal{C}}
\newcommand{\C}{\mathbb{C}}
\newcommand{\Disque}{\mathbb{D}}
\newcommand{\D}{\mathcal{D}}
\newcommand{\Ecal}{\mathcal{E}}
\newcommand{\F}{\mathcal{F}}
\newcommand{\Fcal}{\mathcal{F}}
\newcommand{\Gcal}{\mathcal{G}}
\newcommand{\Hyp}{\mathbb{H}}
\newcommand{\Hil}{\mathcal{H}}
\newcommand{\Hsc}{\textsc{H}}
\newcommand{\Hg}{\textsc{\emph{H}}}
\newcommand{\Ical}{\mathcal{I}}
\newcommand{\K}{\mathbf{K}}
\newcommand{\Kcal}{\mathcal{K}}
\newcommand{\Lin}{\mathscr{L}}
\newcommand{\Lie}{\mathcal{L}}
\newcommand{\Ncal}{\mathcal{N}}
\newcommand{\N}{\mathbb{N}}
\newcommand{\Ocal}{\mathcal{O}}
\newcommand{\Proj}{\mathbb{P}}
\newcommand{\Rcal}{\mathscr{R}}
\newcommand{\R}{\mathbb{R}}
\newcommand{\Scal}{\mathscr{S}}
\newcommand{\Sgot}{\mathfrak{S}}
\newcommand{\Sn}{\mathbb{S}}
\newcommand{\T}{\mathbb{T}}
\newcommand{\Tcal}{\mathcal{T}}
\newcommand{\Ucal}{\mathcal{U}}
\newcommand{\Vcal}{\mathcal{V}}
\newcommand{\Wcal}{\mathcal{W}}
\newcommand{\Z}{\mathbb{Z}}

\newcommand{\dt}{\frac{d}{dt}}
\newcommand{\dd}{\partial}
\newcommand{\SRB}{\mu_{\textsc{SRB}}}

\newcommand{\w}[1]{\langle {#1} \rangle}
\newcommand{\ts}{\textsuperscript}

\DeclareMathOperator{\Cone}{Cone}
\DeclareMathOperator{\dist}{dist}
\DeclareMathOperator{\dv}{div}
\DeclareMathOperator{\Sp}{Sp}
\DeclareMathOperator{\ds}{ds}
\DeclareMathOperator{\Card}{Card}
\DeclareMathOperator{\Ker}{Ker}
\DeclareMathOperator{\Orb}{Orb}
\DeclareMathOperator{\Op}{\emph{Op}}
\DeclareMathOperator{\Opw}{Op}
\DeclareMathOperator{\Int}{Int}
\DeclareMathOperator{\ind}{ind}
\DeclareMathOperator{\Id}{Id}
\DeclareMathOperator{\Graph}{Graph}
\DeclareMathOperator{\Hom}{Hom}
\DeclareMathOperator{\End}{End}
\DeclareMathOperator{\re}{Re}
\DeclareMathOperator{\Res}{Res}
\DeclareMathOperator{\rank}{rank}
\DeclareMathOperator{\Ran}{Ran}
\DeclareMathOperator{\im}{Im}
\DeclareMathOperator{\supp}{supp}
\DeclareMathOperator{\Vect}{Vect}
\DeclareMathOperator{\E}{\boldsymbol{E}}

\def\rest#1#2{\mathchoice
{\setbox1\hbox{${\displaystyle #1}_{\scriptstyle #2}$}
\restrictionaux{#1}{#2}}
{\setbox1\hbox{${\textstyle #1}_{\scriptstyle #2}$}
\restrictionaux{#1}{#2}}
{\setbox1\hbox{${\scriptstyle #1}_{\scriptscriptstyle #2}$}
\restrictionaux{#1}{#2}}
{\setbox1\hbox{${\scriptscriptstyle #1}_{\scriptscriptstyle #2}$}
\restrictionaux{#1}{#2}}}
\def\restrictionaux#1#2{{#1\,\smash{\vrule height.8\ht1 depth.85\dp1}}_{\,#2}}
